\section{Additional new experiment results}
\label{sec:new-results}
These experiments extend the Verified500 comparison in
Table~\ref{tab:verified500} across agent backbones and measure both task
resolution and token use. Earlier campaigns appear in
Appendix~\ref{sec:historical}.

\begin{table}[htbp]
\centering
\caption{New SWE-Context comparison. Entries reproduce the reported
resolved/total counts. The run configuration and benchmark manifest are TBD.}
\label{tab:new-swectx}
\begin{tabular}{lr}
\toprule
Method & Resolved/total \\
\midrule
No memory & 21/98 \\
Free summaries & 19/98 \\
Mem0 & 19/98 \\
LangMem & 24/98 \\
ContextGraph & \textbf{35/98} \\
\bottomrule
\end{tabular}
\end{table}
% Author query: the supplied slide lists35/98 alongside34.69% (+13.26pp).
% These are inconsistent;34.69% corresponds to34/98. Keep the explicitly
% supplied count here until clarified; do not derive a conflicting rate.

\begin{table}[htbp]
\centering
\caption{New Gemini 3.1 Pro/live-SWE-agent comparison. Benchmark identity
and sample count are TBD. The gain column reports relative improvements
from the supplied rates.}
\label{tab:new-gemini}
\begin{tabular}{lrrr}
\toprule
Metric & No memory & ContextGraph & Relative gain \\
\midrule
pass@1 & 82.3\% & \textbf{88.9\%} & +8.0\% \\
pass@3 & 82.6\% & \textbf{90.4\%} & +9.4\% \\
pass@5 & 82.8\% & \textbf{91.7\%} & +10.7\% \\
\bottomrule
\end{tabular}
\end{table}

\paragraph{Token efficiency.}
The new summary reports 236K tokens per task without memory and 201K with
ContextGraph, a 14.8\% reduction (approximately 15\%). The benchmark, sample
count, and accounting boundary for these token totals remain TBD; the source
summary does not identify a specific table as their measurement population.

\begin{table}[htbp]
\centering
\caption{New Codex/GPT-5.5 comparison, preserving the supplied counts and
denominators. Benchmark identity and the pass@3 sampling unit are TBD.}
\label{tab:new-codex55}
\begin{tabular}{lrr}
\toprule
Reported metric & Codex + GPT-5.5 & With ContextGraph \\
\midrule
pass@1 & 1/20 & 3/20 \\
pass@3 & 1/20 & 5/30 \\
\bottomrule
\end{tabular}
\end{table}

\clearpage
\section{Full-context code-agent retrieval}
\label{sec:agentic-retrieval}
The retrieval problem is to find experience that advances the current repair.
A short issue embedding can omit the information that makes a memory useful:
the code already inspected, a failed approach, an observed exception, or a
condition that a partial repair has broken. We therefore propose giving the
selector the coding agent's complete current-task context and tools for
exploring the experience pool.

Let $C_t$ contain the task and conversation, inspected code, edits, actions,
and observations available at step $t$. The selector operates on $C_t$ and
a frozen experience pool $\mathcal{M}$. It can search with FAISS or lexical
queries, open complete source records, follow operation and input relations,
and compare the requirements supplied by different experiences. Its output
is a set of source IDs, their expected contribution to the repair, and the
records or executable conditions supporting that choice. Selection can recur
when new execution feedback changes what the repair needs.

Full context describes the selector's access to the current task; it does
not require placing every experience into one prompt. Search and record
reading make the whole admitted pool accessible. The returned material
retains the source identity and can use the same condition-composition and
execution interface as fixed-candidate retrieval.

\paragraph{Implementation boundary.}
The repository already provides context-aware candidate reranking and a
selector over all 300 compact experience cards. These prototypes cap task
or working-state text at 4,000 characters and limit candidate content.
They provide starting points for the proposed full-context, iterative
code-agent search; the implemented source-witness path currently uses
the FAISS candidate policy described in Section~\ref{sec:method}.

\begin{table}[htbp]
\centering
\caption{\textbf{Does full-context agentic search find more useful experience?}
Planned comparison on the same 24 development tasks, three repetitions per
arm. Freeze a common solver-state snapshot before selection on each task;
keep the model, source pool, delivered memory format, maximum three sources,
composer, and total selection-plus-repair budget fixed. B, D, and E use the
same search-agent controller. B and D can open only the same FAISS ten
candidates; E can search the complete admitted pool. Thus D--B measures
the contribution of full context, and E--D measures iterative pool access.
C compares the existing all-card selector strategy. Evaluate with the common
Combined suite and source-utility protocol of
Tables~\ref{tab:planned-content} and~\ref{tab:planned-utility}; rankings and
selections are frozen before obtaining utility labels. Combined measures
the actual repaired program. Search cost includes all selector calls.}
\label{tab:planned-agentic}
\small
\begin{tabular}{p{0.49\linewidth}ccc}
\toprule
Selection policy & Utility@3 & Combined & Search cost (\$) \\
\midrule
A. FAISS top-three & TBD & TBD & TBD \\
B. Task-only agent; fixed candidates & TBD & TBD & TBD \\
C. Task-only selector; all experience cards & TBD & TBD & TBD \\
D. Full-context agent; fixed candidates & TBD & TBD & TBD \\
E. Full-context agent; iterative pool search & TBD & TBD & TBD \\
\bottomrule
\end{tabular}
\end{table}
